% !PRINTED 86
\weberpage{86}
%<<CONT

zwei Zeilen übereinstimmen, wenn aber $k>\nu$ ist, nach der Voraussetzung, weil dann (12) eine $(\nu+1)$-reihige Determinante der Matrix (3) ist. Wir bekommen also aus (11), da $A$ von Null verschieden ist,
\begin{equation}
 \sum_{\mu=1}^{n}a_\mu^{(k)}x_\mu=0,\qquad k=1,2,\ldots,m,
\tag{13}
\end{equation}
d. h. das System der Gleichungen (1) ist durch (7) befriedigt.

Wir wollen von dem so bewiesenen Satze noch den besonderen Fall hervorheben, dass $m=n-1$ und $\nu=n-1$ ist. In diesem Falle bleibt nur \gesperrt{eine} der Unbekannten beliebig, und die Verhältnisse der Unbekannten sind völlig bestimmt. Wir können diesem Resultate folgenden Ausdruck geben:

Bezeichnen wir die $(n-1)$reihigen Determinanten der Matrix
\begin{equation}
\begin{array}{cccc}
 a_1^{(1)}&a_2^{(1)}&\cdots&a_n^{(1)}\\
 a_1^{(2)}&a_2^{(2)}&\cdots&a_n^{(2)}\\
 \vdots&\vdots&&\vdots\\
 a_1^{(n-1)}&a_2^{(n-1)}&\cdots&a_n^{(n-1)}
\end{array}
\tag{14}
\end{equation}
mit abwechselndem Vorzeichen genommen durch $A_1,A_2,\ldots,A_n$ und nehmen an, dass wenigstens eine von diesen Grössen von Null verschieden sei, so ist die Lösung des Systems:
\begin{equation}
\begin{aligned}
 a_1^{(1)}x_1+a_2^{(1)}x_2+\cdots+a_n^{(1)}x_n&=0,\\
 a_1^{(2)}x_1+a_2^{(2)}x_2+\cdots+a_n^{(2)}x_n&=0,\\
 &\vdots\\
 a_1^{(n-1)}x_1+a_2^{(n-1)}x_2+\cdots+a_n^{(n-1)}x_n&=0
\end{aligned}
\tag{15}
\end{equation}
gegeben durch die Verhältnisse
\begin{equation}
 x_1:x_2:\cdots:x_n=A_1:A_2:\cdots:A_n.
\tag{16}
\end{equation}
So erhalten wir für $n=3$ die Lösung des in der Geometrie oft vorkommenden Gleichungssystems
\begin{equation}
 ax+by+cz=0,\qquad a'x+b'y+c'z=0
\tag{17}
\end{equation}
in der Form
\begin{equation}
 x:y:z=bc'-cb':ca'-ac':ab'-ba'.
\tag{18}
\end{equation}
